% Template LaTeX para TCC de graduação da UNIPAMPA — V2
% Copyright (c) 2026 Charles Quevedo Carpes.
% Autoria das contribuições próprias, organização e adaptações do template.
% This work may be distributed and/or modified under the conditions of the
% LaTeX Project Public License, version 1.3c or (at your option) any later version.
% License: https://www.latex-project.org/lppl/lppl-1-3c.txt (see LICENSE).
% This work has the LPPL maintenance status `author-maintained'.
% The Current Maintainer of this work is Charles Quevedo Carpes.
% Consulte AUTORIA.md para escopo, créditos e relação dos arquivos.

\begin{resumotcc}
Este arquivo deve conter um resumo informativo, redigido em terceira pessoa, em parágrafo único e com extensão entre 150 e 500 palavras. O texto deve apresentar, de maneira concisa, o tema, o objetivo, a metodologia, os principais resultados e as considerações finais do Trabalho de Conclusão de Curso. Evitam-se enumerações, citações, fórmulas, símbolos e abreviaturas pouco conhecidas. Substitua este texto pelo resumo definitivo e confira a correspondência integral entre a versão em português e o resumo em língua estrangeira.

\palavraschave{normalização acadêmica; trabalho de conclusão de curso; Universidade Federal do Pampa; LaTeX}
\end{resumotcc}
